% 第 18 题正方体平面展开图＋直观图（两幅一体）：按固定源内容流矢量层精确值复刻（1 pt → 0.03515 cm，与全库同制）。
% 原点取展开图左缘（横条外轮廓左边）x=174.1018 pt、底缘（直观图 A/B 标签下缘）y（顶向下）=267.1130 pt；y 向上（Y = 267.1130 − y_页）。
% 源实测（figure-inventory §3）：
% 左幅（展开图）：十字六连方（竖条 1×4 叠横条 3×1，共用中格 M）；格边长 28.3466＝1.0000 cm，每格严格正方形；
%   竖条轮廓 x 28.3467–56.6933、y 12.9096–126.3001；横条轮廓 x 0–85.0400、y 69.6068–97.9534；
%   竖条内部第 3 条格线 @ y 41.2602（R3/R4 分界）。8 个字母全部贴网格顶点排印、无一落在格内（§3.1.1）。
% 右幅（直观图）：8 顶点 / 12 棱（9 实 + 3 虚；虚线节距 2.98883，三条均自隐藏顶点 BBL 出发）；
%   前面 68.0374 × 64.1414（宽高比 1.06074 ≠ 1 —— 不是严格斜二测，深度方向却严格 44.995°、
%   长 34.0120 ≈ 前面宽之半；属源内作图约定，照源复刻不改画，见 §7.2 S-1）；
%   M 为 BC 中点（源页 pt (377.7569,244.5318)＝本图 (203.6551,22.5812)），图上另画实心圆点
%   （全图唯一圆点，直径 2.4 pt），字母 M 排于圆点右下。
% 注：除注明「源页 pt」者外，本注释内坐标均为本图坐标系（X = x_页 − 174.1018，Y = 267.1130 − y_页）。
\begin{tikzpicture}[x=0.03515cm,y=0.03515cm,every node/.style={font=\small,inner sep=0pt},
  gl/.style={line width=0.3985bp,line join=round}]
% ===== 左幅：平面展开图（十字六连方，格边长 28.3466）=====
\draw[gl] (28.3467,12.9096) rectangle (56.6933,126.3001);   % 竖条外轮廓（1×4）
\draw[gl] (0,69.6068) rectangle (85.0400,97.9534);          % 横条外轮廓（3×1）
\draw[gl] (28.3467,41.2602)--(56.6933,41.2602);             % 竖条内部第 3 条格线（R3/R4 分界）
% ===== 左幅字母（照源字形盒中心定位；全部锚于格点、不在格内）=====
\node at (20.0097,104.2760) {\(D\)};
\node at (64.1452,104.2760) {\(C\)};
\node at (92.8677,104.2760) {\(G\)};
\node at (-8.1793,61.3570) {\(E\)};
\node at (20.7107,61.3570) {\(A\)};
\node at (64.3752,61.3570) {\(B\)};
\node at (63.7697,40.2960) {\(F\)};
\node at (19.4317,4.6630) {\(H\)};
% ===== 右幅：正方体直观图（顶点照源实测）=====
% 实线 9 条
\draw[gl] (123.5936,10.5572)--(191.6310,10.5572);   % 前面下边（A–B）
\draw[gl] (191.6310,10.5572)--(215.6792,34.6053);   % 右面下边（B–C）
\draw[gl] (123.5936,74.6986)--(191.6310,74.6986);   % 前面上边
\draw[gl] (191.6310,74.6986)--(215.6792,98.7506);   % 右面上边
\draw[gl] (215.6792,98.7506)--(147.6457,98.7506);   % 顶面后边
\draw[gl] (147.6457,98.7506)--(123.5936,74.6986);   % 顶面左边
\draw[gl] (191.6310,74.6986)--(191.6310,10.5572);   % 前/右公共竖棱
\draw[gl] (215.6792,98.7506)--(215.6792,34.6053);   % 后面右边
\draw[gl] (123.5936,74.6986)--(123.5936,10.5572);   % 前面左边
% 虚线 3 条（照源节距 2.98883；均自隐藏顶点 BBL=(147.6457,34.6053) 出发）
\draw[gl,dash pattern=on 2.98883pt off 2.98883pt] (123.5936,10.5572)--(147.6457,34.6053);  % 底面左边（A→隐藏点）
\draw[gl,dash pattern=on 2.98883pt off 2.98883pt] (147.6457,34.6053)--(215.6792,34.6053);  % 底面后边
\draw[gl,dash pattern=on 2.98883pt off 2.98883pt] (147.6457,34.6053)--(147.6457,98.7506);  % 后面左边
% ===== 右幅：BC 中点 M 处的实心圆点（源为字形 • ，中心与 BC 中点逐位吻合）=====
\fill (203.6551,22.5812) circle [radius=1.2000];
% ===== 右幅字母（照源字形盒中心定位）=====
\node at (118.3597,78.6700) {\(E\)};
\node at (140.1347,36.6940) {\(D\)};
\node at (222.2782,33.6440) {\(C\)};
\node at (118.6617,4.6620) {\(A\)};
\node at (196.2972,4.6620) {\(B\)};
\node at (208.0132,16.6880) {\(M\)};
\end{tikzpicture}
